\section{Experiments}
\label{sec:exp}

\subsection{Setup}

Four end-to-end policies: SimLingo (vision--language, $10\times0.25$\,s plans), and the
TransFuser family DiffusionDrive, LTF and DiffusionDriveV2 ($8\times0.5$\,s; V2 additionally
consumes LiDAR). Two further vision--language policies, Alpamayo-R1 and AutoVLA, enter only the
pre-registered rank test of \Cref{sec:negative}; their multi-camera, multi-frame input makes the
occlusion intervention cover $1/12$ to $1/16$ of the visual input rather than all of it, an
asymmetry we correct by masking every camera and history frame in which the object projects,
and report alongside every number derived from them.

Events come from nuScenes and from the OpenScene/nuPlan data underlying NAVSIM, pooled and then
split by \emph{driving side}: left-hand drive is the benchmark side, right-hand drive the
deployment side. The results in \Cref{sec:geometry,sec:behaviour,sec:controls} are established
on the benchmark side ($n{=}312$ lead events); the deployment side holds $8$--$10$ hazard events
and is used only where we say so. \Cref{sec:limits} explains why that number is a property of
the available data rather than of our mining.

\subsection{The three directions form two families}
\label{sec:geometry}

\begin{table}[t]
\centering
\caption{\textbf{Three directions form two families.}
Angles between the hazard-response direction $v_{\mathrm{faith}}$, the braking direction
$v_{\mathrm{decel}}$, and the appearance direction $v_{\mathrm{bright}}$, at each model's
deepest gated layer ($n{=}312$ left-hand-drive lead events). Read every angle against
the split-half floor (what the same data can reproduce) and against $90^\circ$ (two
unrelated directions in high dimensions). $v_{\mathrm{faith}}$ and $v_{\mathrm{decel}}$
cluster; $v_{\mathrm{bright}}$ sits past the random line on the opposite side.}
\label{tab:geometry}
\small
\setlength{\tabcolsep}{4pt}
\begin{tabular}{lccccccc}
\toprule
& & \multicolumn{3}{c}{angle (deg)} & \multicolumn{2}{c}{floor} & \\
\cmidrule(lr){3-5}\cmidrule(lr){6-7}
Policy & $L^*$ & f--d & b--f & b--d & $v_{\mathrm{f}}$ & $v_{\mathrm{b}}$ & rand \\
\midrule
SimLingo & 12 & \textbf{--} & \textbf{133} & \textbf{--} & 9 & 10 & 90 \\
DD & 5 & \textbf{44} & \textbf{96} & \textbf{93} & 5 & 4 & 90 \\
LTF & 6 & \textbf{46} & \textbf{111} & \textbf{93} & 5 & 4 & 90 \\
DDv2 & 6 & \textbf{34} & \textbf{111} & \textbf{115} & 5 & 5 & 90 \\
\bottomrule
\end{tabular}
\end{table}

\Cref{tab:geometry} reports the pairwise angles at each policy's deepest gated layer.
$v_{\mathrm{faith}}$ and $v_{\mathrm{decel}}$ lie $34$--$46^\circ$ apart. Against a split-half
floor of $7^\circ$ this is a real separation, and against the $90^\circ$ random baseline it is a
strong association: the two are near but distinct members of one family. This is consistent
with what we find in \Cref{sec:negative}---that $v_{\mathrm{faith}}$ carries mostly a
\emph{braking-tendency} component and only weakly a \emph{hazard} component.

$v_{\mathrm{bright}}$ behaves differently. Its angle to both action directions is
$93$--$115^\circ$, i.e.\ past the random line, with a mean cosine of $-0.54$ to
$v_{\mathrm{faith}}$ across all gated layers and all four policies (39 of 41 layer-cells
negative). SimLingo is the cleanest case: its cosine sits between $-0.67$ and $-0.76$ across all
24 layers, essentially flat with depth.

The sign is what matters. Since
$v_{\mathrm{faith}} = h(\text{hazard present}) - h(\text{hazard removed})$, a negative
projection says the appearance change moves the representation \emph{toward} the
hazard-removed state. In plain terms, the policy encodes ``it got dark'' the way it encodes
``the hazard is gone.''

\subsection{And the behaviour agrees}
\label{sec:behaviour}

\begin{table}[t]
\centering
\caption{\textbf{Synthetic night makes policies drive \emph{faster}, and the driver is
sensor noise rather than darkness.} $\Delta b$ is the change in planned speed
(m/s, $\arcfull$) when the same frame is re-rendered as night; positive means faster.
``sky'' perturbs only the sky region ($17.9\%$ of the frame, containing no
driving-relevant content); the driving corridor is untouched to the pixel.
$n{=}312$ left-hand-drive lead events.}
\label{tab:deltab}
\small
\begin{tabular}{lcccc}
\toprule
& \multicolumn{2}{c}{sky only} & \multicolumn{2}{c}{whole image} \\
\cmidrule(lr){2-3}\cmidrule(lr){4-5}
Policy & luminance & \;$+$noise & luminance & \;$+$noise \\
\midrule
SimLingo & $-0.126$ & $-0.131$ & $+1.485$ & $+1.572$ \\
DD & $-0.013$ & $+0.045$ & $+0.084$ & $+2.054$ \\
LTF & $-0.103$ & $+0.091$ & $-0.856$ & $+1.245$ \\
DDv2 & $-0.031$ & $-0.020$ & $-0.116$ & $-0.943$ \\
\bottomrule
\end{tabular}
\end{table}

\Cref{tab:deltab} gives the behavioural consequence. Under a full-image night rendering with
noise, DiffusionDrive plans $2.05$\,m/s faster and LTF $1.25$, both with the sign shared by
$>94\%$ of events. For scale:
\emph{deleting the lead vehicle from the image} raises DiffusionDrive's planned speed by
$0.23$\,m/s. The appearance perturbation is therefore roughly nine times more effective at
making this policy accelerate than the removal of the object it is supposed to be reacting to.

Two independent readouts---the negative cosine of \Cref{sec:geometry} and the positive
$\Delta b$ here---agree on direction, in every policy where both are measurable. DiffusionDriveV2
is the informative exception: it is the only LiDAR-consuming policy, and the only one whose
$\Delta b$ is negative ($-0.94$). Its representation is also the most tightly coupled
($\cos = -0.84$ at the output layer). Coupling and behavioural sign are separable; the geometry
says the appearance axis presses on the action axis, not which way.

\paragraph{The mean is representative for only two of the four.} The last two columns of
\Cref{tab:deltab} are there because a mean alone would mislead. For DiffusionDrive the mean
($+2.05$) and median ($+2.12$) coincide, the effect size is $|\mu|/\sigma = 2.12$, and $98\%$
of events share the sign: a concentrated shift. LTF is similar ($1.34$, $94\%$). SimLingo and
DiffusionDriveV2 are not. SimLingo's mean of $+1.57$ sits well above its median of $+0.98$, its
effect size is $0.72$, the top decile of events supplies $31\%$ of the total, and a kernel
density estimate at a bandwidth of $0.4\times$ the reference rule resolves \emph{two} modes, at
$+0.30$ and $+3.98$. The mean falls in the trough between them and describes no event. DDv2 is
weaker still ($|\mu|/\sigma = 0.55$, sign agreement $70\%$).

We read this as the effect being \emph{conditional} in those two policies rather than absent:
the same perturbation leaves a large group of scenes nearly unmoved and drives another group
hard. Which scene property separates the modes---ambient luminance, apparent size of the lead
object, something else---we do not know; identifying it would say more than the aggregate does.
For the claims of this paper we therefore rest on DiffusionDrive and LTF, where the shift is
concentrated, and report SimLingo and DDv2 as directionally consistent but dispersed.

\paragraph{A single event.} In \texttt{scene-1055}, a Singapore night scene, a pedestrian stands
$34$\,m ahead at the edge of oncoming headlight glare; the human brakes from $9.6$ to
$3.8$\,m/s. LTF plans $9.10$\,m/s and DiffusionDrive $8.58$---essentially no deceleration---and
removing the pedestrian from the image changes neither ($b = -0.02$ and $+0.02$). The
representational and behavioural failures we quantify in aggregate are visible in single frames.

\subsection{Noise, not darkness; global, not local}
\label{sec:controls}

Two controls separate the mechanism.

\paragraph{Luminance versus noise.} Columns 1 and 3 of \Cref{tab:deltab} apply the photometric
transform with the noise term set to zero. For DiffusionDrive, sky-only darkening gives
$\Delta b = -0.013$ (marginally \emph{slower}, the safe direction); adding noise flips it to
$+0.045$. Whole-image: $+0.084$ without noise, $+2.054$ with. Darkening is a smooth monotone
map that preserves edge structure up to a scale factor, and a normalised encoder largely undoes
it. Noise adds energy at exactly the spatial frequency where small distant objects live: after
darkening, the frame's own local gradients have $\mathrm{sd}\approx 1.9$, and the injected grain
has $\mathrm{sd}\approx 7$. The evidence is not that darkness is harmless but that the operative
variable is sensor noise---which is precisely what a real high-ISO night camera produces.
SimLingo is again a partial exception: its whole-image response is large ($+1.49$) even without
noise, so for that policy luminance alone suffices.

\paragraph{Sky-only.} The sky mask covers $17.9\%$ of the frame. We verified that the
perturbation stays inside it: high-frequency energy rises from $4.29$ to $5.85$ within the mask
and from $6.33$ to $6.34$ outside, with the with- and without-noise conditions differing by
$10^{-4}$ outside---the residual is feathering, not leakage. The lead vehicle, the road surface
and the lane markings are untouched to the pixel. Yet $\Delta b$ is still $+0.045$ for
DiffusionDrive, $19\%$ of the effect of deleting the hazard entirely. The dependence on
appearance is therefore global---plausibly through a pooled or distributional statistic---rather
than mediated by any cue in the region that carries the driving decision.

\paragraph{Fill colour.} A three-way comparison (grey blob / road colour $+$ noise / road colour
only) on 304 deployment lead events shows the road-colour-only condition reproduces the grey
blob almost digit for digit (DD $+0.234$ vs $+0.214$; LTF $-0.153$ vs $-0.152$), while adding
noise inflates $b$ four- to eight-fold and drives the control arm from $-0.008$ to $+0.204$.
All results here use the noiseless road-colour fill.

\subsection{Projections are not a behavioural proxy}
\label{sec:negative}

\begin{table}[t]
\centering
\caption{\textbf{Pre-registered tests of the projection as a behavioural proxy, and
their outcomes.} Every prediction was written to a timestamped file before the
relevant numbers existed (\texttt{results/PREREG\_*.md}). Four of five are rejected;
the fifth holds for two of four policies. We report these because direction readouts
are routinely used as behavioural proxies without such checks.}
\label{tab:negative}
\footnotesize
\setlength{\tabcolsep}{3pt}
\begin{tabular}{@{}p{0.34\columnwidth}@{\,}c@{\,}p{0.40\columnwidth}@{}}
\toprule
Pre-registered hypothesis & $n$ & Outcome \\
\midrule
Occlusion-scenario angle predicts the official closed-loop rank & 6 policies & \textbf{rejected}: $\rho$ falls from $+1.00$ ($n{=}4$, the set that generated the hypothesis) to $+0.03$ \\
\addlinespace
$v_{\mathrm{faith}}$ separates hazard-caused from red-light hard braking & 30/49 & \textbf{rejected} for 2/4; the other 2 reach $\delta{=}0.28,0.30$ against a pre-set bar of $0.33$ \\
\addlinespace
Representational deviation on the deployment side implies more aggressive behaviour & 48/31 & \textbf{rejected}: 0/4 policies; three are \emph{more} conservative \\
\addlinespace
An axis fit on benchmark hard-braking transfers to deployment & 19/29 $\to$ 11/20 & \textbf{rejected}; the axis fails on its own training split ($\delta{=}0.02$--$0.28$), so no direction exists to transfer \\
\addlinespace
Per-event projection predicts per-event response & 325 & holds for SimLingo ($\rho{=}{+}0.51$) and DDv2 ($+0.25$); null for DD, LTF \\
\bottomrule
\end{tabular}
\end{table}

\Cref{tab:negative} lists five pre-registered hypotheses and their outcomes. The pattern is
consistent: the projection onto $v_{\mathrm{faith}}$ carries some hazard information under
favourable conditions---the lead scenario, large-response policies, quartile-level grouping at
$n{=}325$---and none at the level of individual events, individual policies, or across the
domain shift.

Two of the rejections are worth stating precisely, because the manner of failure is
informative. The rank hypothesis produced a perfect Spearman correlation on the four policies
that generated it and $+0.03$ once Alpamayo-R1 and AutoVLA were added; AutoVLA has the largest
angle and the best closed-loop score, exactly inverted. And the benchmark-to-deployment transfer
failed on its own \emph{training} split (leave-one-out $\delta = 0.02$--$0.28$), so the correct
reading is not that a direction failed to transfer but that at $n{=}19$ versus $29$ in
$256$--$2048$ dimensions there was no direction to begin with.

\begin{table}[t]
\centering
\caption{\textbf{What \emph{does} transfer: a multi-layer speed axis.} Layer weights
and per-layer directions are fit on benchmark cruising samples only; the deployment
side is held out entirely. Spearman $\rho$ between the projection and true ego speed.
The weighted axis beats the best single layer on held-out data, and the advantage
survives the domain shift.}
\label{tab:speed}
\small
\setlength{\tabcolsep}{3.5pt}
\begin{tabular}{@{}lcccccc@{}}
\toprule
& \multicolumn{2}{c}{$n$} & \multicolumn{3}{c}{bench (CV)} & depl. \\
\cmidrule(lr){2-3}\cmidrule(lr){4-6}\cmidrule(lr){7-7}
Policy & b. & d. & best-$L$ & wtd. & shuf. & wtd. \\
\midrule
DD & 390 & 246 & +0.65 & \textbf{+0.69} & +0.11 & \textbf{+0.62} \\
LTF & 390 & 246 & +0.67 & \textbf{+0.72} & +0.20 & \textbf{+0.66} \\
DDv2 & 390 & 246 & +0.57 & \textbf{+0.58} & +0.05 & \textbf{+0.52} \\
\bottomrule
\end{tabular}
\end{table}

\Cref{tab:speed} shows the one readout that does transfer. Fitting per-layer regression
directions against ego speed on benchmark cruising samples, and weighting layers by their
held-out agreement with that target, gives $\rho = +0.58$ to $+0.72$ under 5-fold
cross-validation and $+0.52$ to $+0.66$ on the held-out deployment side, in every case beating
the best single layer and far above a label-shuffled null ($+0.05$ to $+0.20$). The contrast
with the failures above is instructive: this target variable is \emph{external} (measured ego
speed), whereas $v_{\mathrm{faith}}$'s target is the policy's own output.

\subsection{Limits}
\label{sec:limits}

\paragraph{Deployment-side hazard events top out in the low tens.} Four independent criteria
converge: brake-first mining yields 74 right-hand-drive events but none with time-to-collision
below $2.3$\,s; direct TTC screening yields zero below $1.5$\,s once roadside furniture and
cross-traffic are excluded (both verified by inspecting frames); harsh braking
($\ge 3$\,m/s$^2$) yields 31, of which 11 have a visible object ahead; and the entire NAVSIM
right-hand-drive split---1{,}354 scenes---contains no event above $2$\,m/s$^2$. With $8$--$10$
usable events, a bootstrap over the deployment side gives rank probabilities near chance
(27\%/40\%/33\% for one policy across three positions), so we report no deployment-side ranking.

\paragraph{The public sensor release is the bottleneck, not the recording.} Deriving per-city
hours from log durations---nuPlan publishes no city breakdown---the right-hand-drive city
holds $177.3$ hours of recorded logs but only $6.9$ hours with sensors released ($3.9\%$),
against $16.4\%$, $10.0\%$ and $8.8\%$ for the three left-hand-drive cities.

\paragraph{Scope.} \Cref{sec:geometry,sec:behaviour,sec:controls} are established on the
benchmark side. The night transform is closed-form and its noise magnitude is a chosen
parameter, so the magnitude of $\Delta b$ is parameter-dependent; the ordering (noise
$\gg$ luminance) and the sign are not. Whether the effect replicates on the deployment side, and
whether it is calibrated by real night imagery rather than synthesis, are the two open questions
we consider most important.
